% You might need nrc1 instead; but nrc2 should be used usually
\documentclass{nrc2}

\usepackage{graphicx}
\usepackage{amsmath}
\usepackage{amssymb}
\usepackage{bm}
\usepackage{natbib}
\setcitestyle{aysep={},citesep={,}}

\usepackage[english]{babel}

%% Uncomment if your equations need to include the
%% section number, e.g. equation 3.2 is the 2nd equation
%% in section 3.
% \numberby{equation}{section}

%% for nrc1 documentclass
% \renewcommand{\topfraction}{.95}
% \renewcommand{\textfraction}{.05}
% \renewcommand{\floatpagefraction}{.95}

%% for nrc2 documentclass
\renewcommand{\topfraction}{.95}
\renewcommand{\floatpagefraction}{.95}
\renewcommand{\dbltopfraction}{.95}
\renewcommand{\textfraction}{.05}
\renewcommand{\dblfloatpagefraction}{.95}


\title{Title of Your NCR Paper}


%% In this example, Author One, Author Two and Author Four are using th same address
\author{Author One}
\address[label1]{Address One}
\author{Author Two}
%% No address here, so Author Two uses the preceding address i.e. Address One
\author[A.~Three]{Author Three}
\address{Address Three}
%% Author Four uses the address label1
\author{Author Four}
\address[label1]

% Corresponding author's email
\correspond{author@email.address}

\IDnote{Any extra information. Comment out if unnecessary.}

% \author{M. N. Islam}
% \address{Department of Civil \& Environmental Engineering, University of Pittsburgh,Pittsburgh, PA 15213, USA.}
% \author{C. T. Gnanendran}\correspond{xxx@gmail.com}
% \address[label1]{School of Engineering and Information Technology, University of New South Wales,Canberra, ACT 2600, Australia.}
% \author{S. T. Sivakumar}
% \address[label2]{Queensland Department of Transport and Main Roads, Brisbane, QLD 4000, Australia}

\dedication{Dedication, if any. Comment out if unnecessary.}

\abbreviations{abbrvs. if any. Comment out if unnecessary.}

\shortauthor{One et al.}

\usepackage{ragged2e}

\begin{document}
%% Various other fields, such as the journal code and name, will be inserted by the editor after submission.

\maketitle

\begin{abstract}
\justifying
Your abstract goes here. Cupcake ipsum dolor sit amet. Cotton candy lollipop sweet apple pie sugar plum. Halvah jelly-o jujubes tart tart cake sesame snaps. Muffin halvah tootsie roll jelly-o caramels chupa chups. Powder macaroon chupa chups. Lollipop ice cream chocolate cake chupa chups chocolate bar jelly beans gummi bears chocolate jelly-o.



\keywords{keyword 1, keyword 2, keyword 3}

\PACS{1234}
\end{abstract}

\section{Introduction}

Your introduction goes here!

\section{Headings and Subheadings}
\label{sec:examples}

Five levels, numbered automatically.To suppress numbering (e.g., for `Acknowledgements'), use an asterisk before the opening curly brace: \verb|\section*{Acknowledgements}|.

For sub- and superscripts in section titles, use the macros
\verb|\textsubscript| and \verb|\textsuperscript|, respectively.

\begin{enumerate}
\item \verb|\section{...}|
\item \verb|\subsection{...}|
\item \verb|\subsubsection{...}|
\item \verb|\subsubsubsection{...}| (NRC macro)
\item \verb|\paragraph{...}|
\end{enumerate}

\section{Text}

The usual \LaTeX\ commands and environments can be used. Where lists must be flushed to the left margin, there are two NRC-specific environments to use:

\begin{flenumerate} 
\item The \verb|flenumerate| environment creates a numbered list with labels flushed to the left margin.
\item Note that no nesting is possible.
\end{flenumerate}

\begin{flitemize} 
\item The \verb|flitemize| Generates a bulleted list  with labels flushed to the left margin.
\item No nesting possible.
\end{flitemize}

\section{Tables, Figures, Captions}

\subsection{How to Include Figures}

First you have to upload the image file (JPEG, PNG or PDF) from your computer to Overleaf using the upload link the project menu. Then use the includegraphics command to include it in your document. Use the figure environment and the caption command to add a number and a caption to your figure. See the code for Figure \ref{fig:frog} in this section for an example.

\begin{figure}[hbt!]
\centering
%% NRC requires \topcaption to typeset captions at the TOP of floats
\setlength{\captionwidth}{0.3\textwidth}
\topcaption{\label{fig:frog}This frog was uploaded to Overleaf via the project menu.}
\includegraphics[width=0.3\textwidth]{frog.jpg}
\end{figure}

\subsection{How to Make Tables}

Use the table and tabular commands for basic tables --- see Table~\ref{tab:widgets}, for example.

\begin{table}[hbt!]
\centering
%% NRC requires \topcaption to typeset captions at the TOP of floats
\topcaption{\label{tab:widgets}An example table.}
\begin{tabular}{l|r}
\hline*    %% Generates thicker lines -- NRC only.
Item of Interest & Quantity \\\hline
Widgets & 42 \\
Gadgets & 13\\
\hline*    %% Generates thicker lines -- NRC only.
\end{tabular}
\end{table}

\section{How to Write Mathematics}

\LaTeX{} is great at typesetting mathematics. Let $X_1, X_2, \ldots, X_n$ be a sequence of independent and identically distributed random variables with $\text{E}[X_i] = \mu$ and $\text{Var}[X_i] = \sigma^2 < \infty$, and let
$$S_n = \frac{X_1 + X_2 + \cdots + X_n}{n}
      = \frac{1}{n}\sum_{i}^{n} X_i$$
denote their mean. Then as $n$ approaches infinity, the random variables $\sqrt{n}(S_n - \mu)$ converge in distribution to a normal $\mathcal{N}(0, \sigma^2)$.

Equations can be numbered too:

\begin{equation}
S_n = \frac{X_1 + X_2 + \cdots + X_n}{n}
      = \frac{1}{n}\sum_{i}^{n} X_i
\end{equation}

\section{Citations and References}
You can use either

\begin{verbatim}
\begin{thebibliography}
\bibitem ...
\end{thebibliography}
\end{verbatim}
%
or \verb|bibtex|. This template uses \texttt{thebibliography}.

Here's a citation:
\citep{peterman1982}.

\section*{Acknowledgements}
Thank you!

\begin{thebibliography}{}
\bibitem[Peterman(1982)]{peterman1982}
Peterman, R.M. 1982. Model of salmon age structure and its use in  preseason forecasting and studies of marine  survival. Can.~J.~Fish.~Aquat.~Sci.
\textbf{39}: 1444--1452.
\end{thebibliography}

\end{document}